\documentclass[12pt]{article}
\usepackage[utf8]{inputenc}
\usepackage[T1]{fontenc}
\usepackage{hyperref}
\title{Tidore Language Resources for K-12 Curriculum Development}
\author{Compiled by Claude Sonnet 5 (claude-sonnet-5) for Andr\'as Kornai}
\date{}
\begin{document}
\maketitle

\section{Introduction}
Tidore (ISO 639-3 tvo; Glottocode tido1248) is a North Halmahera (West Papuan) language of Maluku Utara, Indonesia. It is closely related to Ternate. Speaker counts range from about 26,000 to about 46,000, and the language is threatened.

\section{Primary Educational Resources}
None found.

\section{Linguistic and Reference Materials}
\begin{itemize}
\item Endangered Languages Project bibliography for Tidore. It lists van Staden's \emph{A Grammar of Tidore} (Mouton, 2010), Voorhoeve's \emph{Struktur Bahasa Tidore} (1994) and Lilipaly et al. (1981): \url{https://endangeredlanguages.com/lang/10935/bibliography}
\item Omniglot page on Tidore: \url{https://omniglot.com/writing/tidore.htm}
\item Two ANU open-research repository documents turned up in the search; their content was not verified: \url{https://openresearch-repository.anu.edu.au/server/api/core/bitstreams/5e57480c-3473-448c-9cca-5f1c6f98b7b7/content} and \url{https://openresearch-repository.anu.edu.au/server/api/core/bitstreams/18aac33f-6ec8-437e-b4ac-830884e1a686/content}
\end{itemize}

\section{Religious and Community Resources}
None found.

\section{Digital and Audio Resources}
Rosetta Panglossia entry: \url{https://rosettapanglossia.longnow.org/index.php?title=Tvo}. No dedicated Tidore dictionary was found.

\section{Curriculum Development Context}
No school materials were found. The grammar by van Staden is the main foundation for a pedagogical description. A dictionary, an orthography and a primer would still be needed.
\end{document}